\newpage
\chapter{ASPECTOS ADMINISTRATIVOS}

\section{Gestión del proyecto}

La gestión del proyecto se realiza bajo la metodología Scrum, con sprints de duración de 2 semanas. El equipo utiliza las siguientes herramientas:

\begin{itemize}
    \item \textbf{Jira:} Para la gestión del backlog, planificación de sprints y seguimiento de tareas.
    \item \textbf{GitHub:} Control de versiones y repositorio central del código.
    \item \textbf{Discord / WhatsApp:} Comunicación diaria del equipo.
    \item \textbf{Google Drive:} Almacenamiento de documentación compartida.
\end{itemize}

El Scrum Master (Hernán) es el responsable de facilitar las ceremonias Scrum y eliminar bloqueos. El equipo se reúne diariamente a las 8:00 am para el Daily Meeting.

\section{Estructura organizacional}

El proyecto cuenta con la siguiente estructura organizacional:

\begin{table}[H]
    \centering
    \caption{Estructura organizacional del proyecto}
    \begin{tabular}{|p{4cm}|p{3.5cm}|p{6.5cm}|}
        \hline
        \textbf{Rol} & \textbf{Responsable} & \textbf{Funciones} \\
        \hline
        Scrum Master & Ramos Alata Hernan Washington (221989) & Facilitar ceremonias Scrum, gestionar Jira, coordinar entregas. \\
        \hline
        Product Owner & Equipo (rotativo) & Definir prioridades del backlog, validar entregables. \\
        \hline
        Desarrollador Backend & Vitorino Marin Efrain (160337) & Base de datos PostgreSQL, APIs, lógica del servidor. \\
        \hline
        Desarrollador Frontend & Mamani Flores Natan (230970) & Interfaces web, UI/UX, conexión con backend. \\
        \hline
        Tester & Polo Chura Marco Rosauro (141632) & Pruebas funcionales, reporte de bugs. \\
        \hline
        Soporte Documentación & Cabeza Huillca Flavio Antony (211265) & Redacción y revisión de documentación técnica. \\
        \hline
    \end{tabular}
\end{table}

\section{Comunicaciones}

\subsection{Canales de comunicación}

\begin{itemize}
    \item \textbf{Daily Meeting (8:00 am):} Reunión diaria de 15 minutos para sincronizar avances y bloqueos.
    \item \textbf{Sprint Planning (lunes de inicio de sprint):} Planificación de tareas del sprint.
    \item \textbf{Sprint Review (viernes final de sprint):} Demostración de avances a los docentes.
    \item \textbf{Sprint Retrospective (viernes final de sprint):} Mejora continua del proceso.
\end{itemize}

\subsection{Recursos técnicos y despliegue}

\textbf{Estructura del repositorio:}
\begin{verbatim}
tupaunsaac/
|-- backend/
|   |-- .env.example
|   |-- database/
|   |   |-- schema.sql
|   |   |-- migrate.ts
|   |   `-- seed.ts
|   `-- src/
|-- frontend/
|   |-- .env.example
|   `-- src/
|-- docs/
|   `-- Plan_Proyecto.pdf
`-- docker-compose.yml
\end{verbatim}

\textbf{Variables de entorno:}
\begin{verbatim}
# Backend .env
NODEENV=development
BACKENDURL=http://localhost:8080
FRONTENDURL=http://localhost:3000
PORT=8080

DBDIALECT=postgres
DBHOST=localhost
DBPORT=5432
DBUSER=tupa
DBPASS=tupa
DBNAME=tupa

JWTSECRET=cambia-este-secreto
REDISURI=redis://:tupa@127.0.0.1:6379

# Frontend .env
VITEAPIBASEURL=http://localhost:8080
VITEAPPNAME=Portal de Tramite Documentario UNSAAC
\end{verbatim}

\subsection{Configuración y despliegue}

\textbf{Desarrollo local:}
\begin{verbatim}
docker compose up -d postgres redis
cd backend && npm install && npm run dev
cd frontend && npm install && npm run dev
\end{verbatim}

\textbf{Base de datos:}
\begin{verbatim}
cd backend
npm run db:migrate   # Ejecuta schema.sql
npm run db:seed      # Carga datos de prueba
npm run db:setup     # Migrate + Seed
\end{verbatim}

\textbf{Despliegue en producción (Ubuntu 24.x):}
\begin{verbatim}
# Instalar dependencias
sudo apt-get update && sudo apt-get upgrade -y
sudo apt-get install -y build-essential postgresql-client git

# Instalar Node 24.x
curl -fsSL https://deb.nodesource.com/setup_24.x | sudo -E bash -
sudo apt-get install -y nodejs

# Instalar Docker
curl -fsSL https://get.docker.com -o get-docker.sh
sudo sh get-docker.sh
sudo usermod -aG docker ${USER}

# PostgreSQL con Docker
docker volume create tupa_postgres_data
docker run --name tupa-postgres -e POSTGRES_DB=tupa -e POSTGRES_USER=tupa -e POSTGRES_PASSWORD=tupa -p 5432:5432 -d --restart=always -v tupa_postgres_data:/var/lib/postgresql/data postgres:17

# Redis con Docker
docker run --name tupa-redis -p 6379:6379 -d --restart=always redis:latest redis-server --appendonly yes --requirepass "tupa"

# Clonar y desplegar
git clone https://github.com/lovito99/tupaunsaac.git
cd tupaunsaac/backend && npm install && npm run build && npm run db:migrate
cd ../frontend && npm install && npm run build
\end{verbatim}

\subsection{Enlaces útiles}

\begin{itemize}
    \item \textbf{Repositorio GitHub:} \url{https://github.com/Natanmamani/ProyectoDesarrollo}
    \item \textbf{Prototipo Figma:} \url{https://figma.com/...} (Agregar enlace real)
    \item \textbf{Jira:} \url{https://...} (Agregar enlace real)
\end{itemize}

\subsection{Contacto de soporte}

Para consultas relacionadas con el despliegue y soporte técnico, contactar a: \texttt{lovito99\_m@live.com}
